\documentclass[10pt,a4paper]{amsart}
\usepackage[margin=19mm]{geometry}
\usepackage{amsmath,amssymb,mathtools}
\usepackage[scr=rsfs,cal=euler]{mathalfa}
\usepackage{xfrac}
\usepackage[unicode,hidelinks,bookmarks=false]{hyperref}
\newcommand{\ie}{i.e.\ }
\newcommand{\cf}{cf.\ }
\newcommand{\resp}{respectively\ }
\newcommand{\Subsection}{\subsection}
\providecommand{\nobreakdash}{\nobreak\hbox{-}}
\setlength{\emergencystretch}{2em}
\multlinegap0pt
\usepackage{amssymb}
\usepackage{enumerate}
\usepackage{xcolor}
\newtheorem{lemma}{Lemma}
\newtheorem{theorem}{Theorem}
\newcommand{\A}{ {\hat A}}
\newcommand{\Amo}{{\mathcal {\bf A}}}
\newcommand{\Ro}{{\overline R}}
\newcommand{\So}{{\hat {\mathcal F}}}
\newcommand{\Vo}{{\bf V}}
\newcommand{\aR}{{\hat R}}
\newcommand{\aV}{{\hat V}}
\newcommand{\ad}{{\rm ad}}
\newcommand{\ba}{{\bf a}}
\newcommand{\bv}{{\bf v}}
\newcommand{\hJ}{{\hat J}}
\newcommand{\hW}{{\hat W}}
\newcommand{\ha}{{\hat a}}
\newcommand{\hv}{{\hat v}}
\newcommand{\hw}{{\hat w}}
\title{$2$-generated axial algebras of Monster type}
\author{Clara Franchi, Mario Mainardis, Sergey Shpectorov}
\date{Source: arXiv:2101.10315v3 (CC BY 4.0)}
\begin{document}
\maketitle

\noindent\emph{Source and licence.} $2$-generated axial algebras of Monster type. Version arXiv:2101.10315v3 (CC BY 4.0), distributed by arXiv under the Creative Commons Attribution 4.0 International licence, http://creativecommons.org/licenses/by/4.0/, as recorded in the arXiv licence metadata for this exact version and shown on the abstract page https://arxiv.org/abs/2101.10315v3. Full licence notice: \texttt{licenses/arxiv-2101.10315-CC-BY-4.0.txt}.\par\smallskip

\noindent\textit{Editorial scope.} Counted unit: the article's theorem free (source0.tex lines 845--847). The article's complete original proof of it is delivered with it (source0.tex lines 849--892, one \texttt{proof} environment of 44 lines). No mathematics is written, repaired or re-derived: the statement and the proof are the article's own lines, in the article's own notation, and no retained line is substituted or re-flowed.  Retained uncounted as the source-local context the proof needs: lines 476--490; lines 514--528; lines 735--737.  Nothing was omitted from any retained span, so the package carries no omission record.  No retained line contains a citation, so the wrapper carries no bibliography and no bibliography entry.  No retained line contains a figure environment, an image-inclusion command or drawing code, so no image is delivered and the package declares no asset dependency.  Editorial formatting only, in addition: an AMS \texttt{amsart} preamble that restates the article's own declarations and macros used by the retained lines, namely the environments \texttt{lemma}, \texttt{theorem} on separate counters, together with the standard packages those lines need.  The retained environments print in this document as Lemma 1, Theorem 1 in order of appearance, whereas the article prints the same items with the numbering of its own full document.  The complete native source of the article as distributed in the arXiv e-print for this exact version is bundled unchanged as \texttt{sources/107371/source0.tex}.
\begin{lemma}\label{lambdafunction}
Let $(R,V, A, (\mathcal F, \star))$ be a primitive axial algebra and assume that the elements of $\mathcal F$ are pairwise distinguishable and the axes in $A$ are free. Then, for every $a\in A$, there is a well defined $R$-linear map 
\begin{equation}
\label{lambda}
\begin{array}{rcccc}
\lambda_a & \colon & V& \to & R\\
&& v &\mapsto & \lambda_a(v)
\end{array}
\end{equation}
such that every $v\in V$ decomposes uniquely as
\begin{equation}\label{justapply}
v=\lambda_a(v)a+\sum_{\mu\in \mathcal F\setminus \{1\}}v_\mu, 
\end{equation}
with $v_\mu\in V_\mu^a$.
\end{lemma}

\begin{lemma}\label{condition}
Let $V$ be an algebra over a ring $R$ and $ (\mathcal F, \star)$ a fusion law with $\mathcal F\subseteq R$ and let $a$ be an idempotent in $V$.  Assume that the elements of $\mathcal F$ are pairwise distinguishable, $Ann_R(a)=\{0_R\}$, and $\ad_a$ is semisimple with spectrum $\mathcal F$. Then, 
\begin{enumerate}
\item $a$ is an $\mathcal F$-axis if and only if, for every  $\gamma, \delta\in \mathcal F$, and $v, w\in V$ 
\begin{equation}\label{cond2}
 (\prod_{\eta \in \gamma\star \delta}(\ad_a-\eta )) (f_\gamma(\ad_a)( v) \cdot f_\delta(\ad_a) (w))=0.
\end{equation}
%for every  $\gamma, \delta\in \mathcal F$, and $v, w\in V$;\\
\item if $a$ is an axis, then $a$ is primitive if and only if, for every  $w\in V$, there exists $l_w\in R$ such that
 \begin{equation}\label{cond1}
 f_1(\ad_a)(w-l_w a) =0.
 \end{equation}
In particular, if  $a$ is primitive, $l_w=\lambda_w(a)$.
\end{enumerate}
\end{lemma}

\begin{lemma}\label{contained}
Let ${\mathcal V}:=(R, V, A, (\mathcal F, \ast) )$ be an element of  $\mathcal O_L$, let $\phi\colon \aR\cup\aV  \to R\cup V$ be a map  that satisfies conditions (H1) (with respect to $(\So,\star)$ and $(\mathcal F, \ast)$), (H2), and (H3). Then $I_0\subseteq \ker \phi_{|_{\aR}}$,  $\hJ_0\subseteq  \ker \phi_{|_{\aV}}$. 
\end{lemma}

\begin{theorem}\label{free}
$\mathcal V_L$ is an initial object in the category $({\mathcal O}_L,{\mathcal Mor}_L)$.
\end{theorem}

\begin{proof}
Assume  ${\mathcal  V}:=(R, V, A, (\mathcal F, \ast))$ is an object in $\mathcal O_L$. By condition O4 of the definition of $\mathcal O_L$, there is a unitary ring homomorphism  $\zeta:\hat D\to R$  inducing a bijection between $\{\hat x_1,  \ldots , \hat x_n\}$ and $\mathcal F\setminus \{0,1\}$. Let $A:=(b_1, \ldots , b_k)$ and let $t$ be the map $\ha_i\mapsto b_i$ from $\A$ to $A$ and 
$
\bar t$ be the map $\ba_i\mapsto b_i$ from $\Amo$ to $\A$. Since $\hW$ is the free commutative magma generated by $\A$, there is a unique magma homomorphism 
$$\chi:\hW\to V,$$ extending the map $t$. Since the elements of $\Lambda$ are algebraically independent over $\hat D$, there is a unique homomorphism of $\hat D$-algebras 
 $$\hat \psi :\aR\to R$$ extending $\zeta$ and 
  such that,  for $i\in \{1,\ldots , k\}$ and  $\hw \in \hW\setminus \{\ha_i\}$,  
  \begin{equation}\label{lambdax}
 \lambda_{\ha_i, \hw}^{\hat \psi}=\lambda_{b_i}(\hw^\chi),
 \end{equation}
  where $\lambda_{b_i}$ is the function defined in Lemma~\ref{lambdafunction}.  Define 
$$
\begin{array}{rcccc}
\hat \chi &:& \aV& \to &V\\
& & \sum_{\hw\in \hW}\gamma_\hw \hw & \mapsto & \sum_{\hw\in \hW}\gamma_\hw^{\hat \psi}\: \hw^\chi
\end{array}.
$$
Then $\hat \chi$ is a ring homomorphism extending $t$. Moreover,  for every $\gamma\in \aR$ and $\hv=\sum_{\hw\in \hW}\gamma_\hw \hw\in \aV$, we have 
\begin{eqnarray*}
(\gamma  \hv)^{\hat \chi}&=&(\gamma \sum_{\hw\in \hW} \gamma_\hw \hw)^{\hat\chi}=( \sum_{\hw\in \hW} \gamma \gamma_\hw \hw)^{\hat\chi}\\
&=&\sum_{\hw\in \hW}\gamma^{\hat \psi} \gamma_\hw^{\hat \psi} \hw^\chi=\gamma^{\hat \psi}(\sum_{\hw\in \hW} \gamma_\hw^{\hat \psi} \hw^\chi)=\gamma^{\hat \psi} v^{\hat \chi}.\\
\end{eqnarray*}
Since, for every $\ha_i \in \A$, $\hv\in \aV$, and $\gamma \in \hat {\mathcal F}$,
$$
 [ (\ad_{\ha_i}-\gamma\: {\rm id}_\aV) (\hv)]^{\hat \chi} =(\ad_{b_i}-\gamma^{\hat \psi}\: {\rm id}_{V}) (\hv^{\hat \chi}),
$$
by Lemma~\ref{condition}, it follows that $\hJ_0$ is contained in $\ker \hat \chi$. Hence $\hat \chi$ induces a ring homomorphism 
$$\begin{array}{rccc}
\phi_V : &\Vo& \to &V\\
& \hv+\hJ_0 & \mapsto & \hv^{\hat \chi}
\end{array}
$$
 extending $\overline t$. Moreover,  for every $\gamma\in \aR$ and ${\bv}\in \Vo$,  we have
$$
(\gamma \bv)^{\phi_{V}}=(\gamma \hv)^{\hat\chi}=\gamma^{\hat \psi} \hv^{\hat \chi}=\gamma^{\hat \psi} {\bv}^{\phi_{V}}. 
$$
From this equality, as in the proof of Lemma~\ref{contained}, we get that $I_0\subseteq \ker \hat \psi$. Then  $\hat \psi$ induces a well defined homomorphism of $\hat D$-algebras $\phi_{R}:\Ro\to R$, $r+I_0\mapsto r^{\hat\psi}$. It follows that  the map $\phi\colon \Ro\cup \Vo \to R\cup V$ such that $\phi_{|\Ro}=\phi_R$ and $\phi_{|\Vo}=\phi_V$ belongs to ${\mathcal Mor}({\mathcal V}_L, \mathcal V)$ . 

Since $\aR=\hat D[\Lambda]$, $\phi_R$ is completely determined by its values on the elements $\lambda_{\ha,\hw}+I_0$, with $\ha\in \A$ and $\hw\in \hW\setminus \{\ha\}$. Further, by Equation~(\ref{lambdax}), for every $\ha_i\in \A$, $\hw\in \hW\setminus \{\ha_i\}$, 
\begin{equation}\label{bendef}
(\lambda_{\ha_i,\hw}+I_0)^{\phi_R}=\lambda_{\ha_i,\hw}^{\hat\psi}=\lambda_{b_i}(\hw^\chi)=\lambda_{b_i}((\hw+\hJ_0)^{\phi_{V}}),
\end{equation}
whence  $\phi_R$ is completely determined by the images $(\hw+\hJ_0)^{\phi_{V}}$, with $\hw \in \hW$.  Since $\phi_{V}$ is a ring homomorphism extending $t$, such images are uniquely determined by $t$, whence the uniqueness of  $\phi$. 
\end{proof}

\end{document}
